\documentclass[UTF8,fontset=fandol,11pt]{ctexart}
\usepackage[a4paper,margin=24mm]{geometry}
\usepackage{amsmath,amssymb,amsthm,graphicx,tabularx,booktabs,hyperref}
\newtheorem{definition}{定义}[section]
\newtheorem{theorem}[definition]{定理}
\newtheorem{example}[definition]{例题}
\newtheorem{exercise}[definition]{练习}
\hypersetup{unicode=true,colorlinks=true,linkcolor=blue}
\title{Offline draft recovery}
\date{}
\begin{document}
\maketitle
\par\noindent 离线合成演示：不代表真实模型对本专题的教学验收。\par
\section{Definition}\label{sec:S01}
% UNIT S01.U00
\begin{definition}
正整数集上的实值函数称为实数列。
\end{definition}
% UNIT S01.U05
\begin{example}
判断交错调和级数
\[
\sum_{n=1}^{\infty}\frac{(-1)^{n-1}}{n}
\]
的敛散性，并估计其前 $n$ 项部分和的绝对误差。

作为通项取值的单点核对，$n=3$ 时有
\[
\frac{(-1)^{3-1}}{3}=\frac13.
\]
这一数值核对只验证第三项的取值，不能证明级数收敛，也不能证明一般的余项估计。收敛性与误差界须由交错级数判别法及其余项估计在核实条件后得到，具体论证见本例解答。
\end{example}
% UNIT S01.U06
\par\noindent\textbf{解：}令 $b_n=1/n$。对每个正整数 $n$，有 $b_n>0$，且
\[
b_{n+1}=\frac{1}{n+1}<\frac1n=b_n,
\qquad \lim_{n\to\infty}b_n=0.
\]
因此由交错级数判别法，级数
\[
\sum_{n=1}^{\infty}\frac{(-1)^{n-1}}{n}
\]
收敛。若其和为 $S$，第 $n$ 个部分和为 $S_n$，则由交错余项估计
\[
|S-S_n|\leq b_{n+1}=\frac{1}{n+1}.
\]
例如，取 $n=99$ 时，前 $99$ 项部分和与级数和的绝对误差不超过 $1/100$。

还需区分收敛与绝对收敛：若级数各项的绝对值所成级数收敛，则称原级数绝对收敛；若原级数收敛而其绝对值级数发散，则称原级数条件收敛。对于本例，
\[
\sum_{n=1}^{\infty}\left|\frac{(-1)^{n-1}}{n}\right|=\sum_{n=1}^{\infty}\frac1n
\]
是发散的调和级数。因此该交错调和级数收敛但不绝对收敛，属于条件收敛。这里判别法说明收敛性，余项估计说明截断误差的大小。\par
\end{document}
